\subsection{Discussion}

The results of our experiments with the PaperSwarm system reveal that ridge regression outperforms the minimum-norm ordinary least squares (OLS) solution in terms of test mean squared error (MSE). Specifically, the relative reduction in mean test MSE achieved by ridge regression over OLS is \result{cl-improve} percent, as reported in \citet{fars2026}. This improvement is attributed to the ability of ridge regression to mitigate the variance in the OLS solution, which is subject to high variability as indicated by the standard deviation of the OLS test MSE across seeds, \result{cl-ols-std}.

The selected ridge penalty, \result{cl-alpha}, was chosen to minimize the mean test MSE, demonstrating the system's capability to adaptively tune hyperparameters. The low standard deviation of the ridge test MSE across seeds, \result{cl-ridge-std}, indicates that the selected penalty provides a more stable and reliable solution compared to OLS.

The design matrix in our experiments contained \result{cl-features} features, and the system was tested with \result{cl-train} training samples per split and \result{cl-test} test samples per split. These configurations were averaged over \result{cl-seeds} random seeds, ensuring robustness and generalizability of the results. The irreducible noise floor, \result{cl-noise}, serves as a reference for the minimum achievable MSE, highlighting the effective reduction in noise achieved by the ridge regression solution.

In conclusion, our experiments with PaperSwarm demonstrate the potential of evidence-gated multi-agent systems to generate reproducible and reliable research. The system's ability to adaptively tune hyperparameters and reduce variance in the solution is a significant step towards more robust and generalizable machine learning models. Future work could explore the application of similar techniques to other machine learning tasks and datasets, further validating the effectiveness of our approach.
